\documentclass[11pt]{article}
\usepackage[margin=1in]{geometry}
\usepackage{amsmath,amssymb,amsthm}
\usepackage{algorithm}
\usepackage{algpseudocode}
\usepackage{booktabs}
\usepackage{hyperref}

\newtheorem{definition}{Definition}
\newtheorem{theorem}{Theorem}
\newtheorem{lemma}[theorem]{Lemma}
\newtheorem{property}{Property}

\title{\textbf{Mathematical Specification, Space-Time Pareto Frontiers, and Segment Recomputation of Gradient Checkpointing (ALGO-NN-27)}}
\author{Team Scaibu \\ \small Email: \texttt{jaydeep.9664920749@gmail.com} \\ \small Architecture and Core Engine Team}
\date{\today}

\begin{document}
\maketitle

\begin{abstract}
Standard backpropagation stores all $L$ intermediate layer activations, consuming $\mathcal{O}(L)$ peak GPU memory. Gradient Checkpointing stores only $\mathcal{O}(\sqrt{L})$ activations at segment boundaries and recomputes intermediate values on-demand during the backward pass at the cost of one additional forward pass ($+33\%$ compute overhead). We present the formal space-time optimization derivation, prove the $\sqrt{L}$ Sublinear Memory Theorem, provide full pseudocode matching the reference implementation, and derive complexity.
\end{abstract}

\section{Notation and Mathematical Preliminaries}
\label{sec:notation}

Let $L \in \mathbb{N}$ denote the total number of sequential computational layers $f_1, \dots, f_L$ in a deep architecture, with initial activation $x_0 \in \mathbb{R}^D$ and recurrence $x_l = f_l(x_{l-1})$.

\subsection{Checkpoint Boundary Mechanics}
\paragraph{Intuitive Meaning:} Instead of storing every layer output in memory, we save activations only at selected checkpoints spaced $k = \lceil\sqrt{L}\rceil$ layers apart. When backpropagation reaches a segment, it recomputes that segment's activations from the saved checkpoint and immediately frees them after use.

\paragraph{Formal Mathematical Definition:}
Let $k \in \mathbb{N}$ denote segment length. Checkpoints are stored at layer indices $\mathcal{C} = \{l \in \{0, \dots, L-1\} \mid l \equiv 0 \pmod k\}$.
The total stored checkpoint activations count is $N_c = \lceil L / k \rceil$.
During backward evaluation of any segment $[i \cdot k, (i+1) \cdot k - 1]$, at most $k$ intermediate activations are recomputed and kept in working memory.
Total peak memory is:
\begin{equation}
\label{eq:mem_objective}
M(k) = \frac{L}{k} + k.
\end{equation}

\paragraph{Worked Numerical Example:}
Let $L = 100$.
If $k = 10$: $N_c = 10$, segment recomputation memory $= 10$. Peak memory $= 10 + 10 = 20$ activations (an $80\%$ reduction compared to full caching of 100 activations).

\subsection{Summary of Symbols}
\begin{itemize}
    \item $L \in \mathbb{N}$: Total number of sequential layers.
    \item $k \in \mathbb{N}$: Segment step interval ($\sqrt{L}$).
    \item $x_0 \in \mathbb{R}$: Input initial activation.
    \item $\mathcal{C} \subset \{0, \dots, L-1\}$: Saved checkpoint index set.
    \item $\mathbf{x}_{\mathrm{check}} \in \mathbb{R}^{N_c}$: Stored checkpoint values.
    \item $x_L \in \mathbb{R}$: Final network output.
    \item $R_{\mathrm{saved}} \in [0, 1]$: Theoretical memory reduction ratio $1 - \frac{\sqrt{L}}{L}$.
\end{itemize}

\section{Problem Definition and Core Concepts}
\label{sec:problem_definition}

\subsection{Optimal Segment Size Optimization}
Minimizing total peak memory $M(k) = \frac{L}{k} + k$ with respect to continuous relaxation $k \in \mathbb{R}_{>0}$:
\begin{equation}
\frac{d M}{d k} = -\frac{L}{k^2} + 1 = 0 \implies k^* = \sqrt{L}.
\end{equation}
Substituting $k^* = \sqrt{L}$ into $M(k)$ yields optimal peak memory $M^* = \sqrt{L} + \sqrt{L} = 2\sqrt{L} = \mathcal{O}(\sqrt{L})$.

\section{Algorithmic Specification}
\label{sec:algorithm}

Algorithm~\ref{alg:grad_checkpoint} specifies the forward checkpointing routine matching \texttt{impl.py}.

\begin{algorithm}[H]
\caption{Gradient Checkpointing Forward Execution and Checkpoint Stash}
\label{alg:grad_checkpoint}
\begin{algorithmic}[1]
\Require Layer count $L \ge 1$, initial activation $x_0$, optional segment size $k_{\mathrm{in}}$.
\Ensure Checkpoint values $\mathbf{x}_{\mathrm{check}}$, final output $x_L$, checkpoint indices $\mathcal{I}_{\mathrm{check}}$, memory saved ratio $R_{\mathrm{saved}}$.
\If{$k_{\mathrm{in}} > 0$}
    \State $k \gets k_{\mathrm{in}}$
\Else
    \State $k \gets \max(1, \lfloor\sqrt{L}\rfloor)$
\EndIf
\State Initialize $\mathbf{x}_{\mathrm{check}} \gets [], \mathcal{I}_{\mathrm{check}} \gets []$
\State $curr \gets x_0$
\For{$l = 0$ to $L-1$}
    \If{$l \bmod k = 0$}
        \State Append $curr$ to $\mathbf{x}_{\mathrm{check}}$, append $l$ to $\mathcal{I}_{\mathrm{check}}$
    \EndIf
    \State $curr \gets f_l(curr)$ \Comment{Evaluate layer $l$}
\EndFor
\State $R_{\mathrm{saved}} \gets \max\left(0.0, 1.0 - \frac{|\mathbf{x}_{\mathrm{check}}|}{L}\right)$
\State \Return $\mathbf{x}_{\mathrm{check}}, \quad curr, \quad \mathcal{I}_{\mathrm{check}}, \quad R_{\mathrm{saved}}$
\end{algorithmic}
\end{algorithm}

\section{Theoretical Guarantees and Memory Bounds}
\label{sec:guarantees}

\begin{theorem}[Sublinear Memory Scaling Theorem]
\label{thm:sublinear_memory}
\textbf{Assumptions:} Let $L \ge 1$ homogeneous sequential layers require unit memory per forward activation.
\textbf{Guarantee:} Gradient Checkpointing with segment length $k = \lceil\sqrt{L}\rceil$ evaluates exact gradients using $\mathcal{O}(\sqrt{L})$ peak activation memory at the cost of exactly one additional forward pass ($+33\%$ total computation time).
\textbf{Proof:}
Storing boundary checkpoints requires $\lceil L / \sqrt{L} \rceil = \mathcal{O}(\sqrt{L})$ memory. During backward propagation, recomputing segment activations from checkpoint $x_{i \cdot k}$ to $x_{(i+1) \cdot k}$ requires $k = \mathcal{O}(\sqrt{L})$ working memory. Because intermediate segment activations are deallocated immediately upon completing backward steps through that segment, peak memory is $\mathcal{O}(\sqrt{L}) + \mathcal{O}(\sqrt{L}) = \mathcal{O}(\sqrt{L})$.
Each layer is computed once in the global forward pass and once in the segment recomputation pass, adding $L$ forward layer evaluations. Since standard backward pass costs $\approx 2L$ forward operations, total computation changes from $1L + 2L = 3L$ to $1L + 1L + 2L = 4L$, representing a $\frac{4L - 3L}{3L} = +33.3\%$ compute overhead.
\textbf{Limitations:} Non-deterministic operations (e.g., Dropout) must preserve identical random seeds across forward and recomputation passes.
\end{theorem}

\section{Complexity Analysis}
\label{sec:complexity}

\begin{itemize}
    \item \textbf{Time Complexity:} $\mathcal{O}(L)$ layer evaluations forward, $\mathcal{O}(L)$ layer recomputations backward. Total time is $\mathcal{O}(L)$ with $1.33\times$ constant factor.
    \item \textbf{Space Complexity:} $\mathcal{O}(\sqrt{L})$ peak memory.
\end{itemize}

\section{Worked Numerical Example}
\label{sec:worked_example}

Let $L = 4, x_0 = 1.0, k = \lfloor\sqrt{4}\rfloor = 2$, layer operation $f(x) = \tanh(0.9x + 0.1)$:
\begin{enumerate}
    \item $l = 0: 0 \bmod 2 = 0 \implies \text{Checkpoint 0 saved at } l=0 \text{ with } x = 1.0$.
    $x_1 = \tanh(0.9(1.0) + 0.1) = \tanh(1.0) = 0.761594$.
    \item $l = 1: 1 \bmod 2 = 1 \implies \text{No checkpoint}$.
    $x_2 = \tanh(0.9(0.761594) + 0.1) = \tanh(0.785435) = 0.655993$.
    \item $l = 2: 2 \bmod 2 = 0 \implies \text{Checkpoint 1 saved at } l=2 \text{ with } x = 0.655993$.
    $x_3 = \tanh(0.9(0.655993) + 0.1) = \tanh(0.690394) = 0.598285$.
    \item $l = 3: 3 \bmod 2 = 1 \implies \text{No checkpoint}$.
    $x_4 = \tanh(0.9(0.598285) + 0.1) = \tanh(0.638457) = 0.563870$.
    \item Checkpoints saved: 2 out of 4 $\implies R_{\mathrm{saved}} = 1 - \frac{2}{4} = 50\%$.
\end{enumerate}

\section{Numerical Considerations and Edge Cases}
\label{sec:numerical}

\begin{itemize}
    \item \textbf{Small Depth Guard:} For $L \le 3$, $k = 1$ saves every layer, gracefully falling back to standard full caching without branch divergence.
\end{itemize}

\begin{thebibliography}{9}
\bibitem{chen2016} T.~Chen et al., ``Training Deep Nets with Sublinear Memory Cost,'' \emph{arXiv:1604.06174}, 2016.
\end{thebibliography}

\end{document}
